% Report skeleton for Linear and Logistic Regression, WS 2026/27.
\documentclass[11pt,a4paper]{article}
\usepackage[utf8]{inputenc}
\usepackage[T1]{fontenc}
\usepackage[margin=1in]{geometry}
\usepackage{amsmath}
\usepackage{amssymb}
\usepackage{amsfonts}
\usepackage{graphicx}
\usepackage{enumitem}
\usepackage{booktabs}
\usepackage{graphicx}
\usepackage{bm}

\newcommand{\sheetnumber}{2}
\newcommand{\R}{\mathbb{R}}

\usepackage{listings}

\providecommand{\xx}{\mathbf{x}}
\providecommand{\yy}{\mathbf{y}}
\providecommand{\XX}{\mathbf{X}}
\providecommand{\tth}{\boldsymbol{\theta}}
\providecommand{\sigm}{\sigma}
\providecommand{\rank}{\operatorname{rank}}

\lstset{
  basicstyle=\ttfamily\small,
  frame=single,
  breaklines=true,
  columns=fullflexible
}

\begin{document}

\begin{titlepage}
  \begin{center}
    {\Large \textbf{Exercise Sheet \sheetnumber}}

    \vspace{0.5em}
    Machine Learning for AIE, WS 2026/27

    \vspace{2em}
    {\LARGE \textbf{Linear and Logistic Regression}}
  \end{center}

  \vspace{2em}
  \begin{center}
    \begin{tabular}{|l|l|l|}\hline
      \multicolumn{3}{|c|}{\textbf{Team Members}} \\\hline
      Last name & First name & Matriculation number \\\hline
      ZENZ & Sebastian & 01031159 \\\hline
    \end{tabular}
  \end{center}
\end{titlepage}

\tableofcontents
\clearpage

\section{Task 1: Linear Regression- Closed-Form Solution}
\label{sec:closedform}

\subsection{Task 1.1: Pen and paper computation}
\label{subsec:closedform1}
For task 1.1 a toy data set of $N = 4$ houses is examined. The input $x$ is the number of rooms and the target $y$ is the price (in 100\,k\$):
\begin{itemize}[itemsep=0.1em]
  \item Task 1.1(a) The design matrix $\XX$ and target variable $\yy$ are defined. $\XX^T\XX$, $\XX^T\yy$, $\tth^* = (\theta_0^*, \theta_1^*)^T$, $\hat{y}^{(i)} = \tth^{*T}\xx^{(i)}$ and the loss $\mathcal{L}(\tth^*)$ are computed. 
  \item Task 1.1(b) A second feature $x_2 = 2x_1$ (number of windows) was added to the design matrix. 
      The rank of this adjusted design matrix $\widetilde{\mathbf{X}}$ was computed and the invertability was discussed. 
      Also three different parameter sets $\tth \in \mathbb{R}^3$ were found, that reach the minimal loss from (a)
\end{itemize}

\begin{center}
\begin{tabular}{c@{\hspace{2em}}cc}
    \toprule
    house $i$ & input $x^{(i)}$ (rooms) & target $y^{(i)}$ (price) \\
    \midrule
    1 & 1 & 1 \\
    2 & 2 & 2 \\
    3 & 3 & 2 \\
    4 & 4 & 3 \\
    \bottomrule
\end{tabular}
\end{center}

\begin{figure}[h]
  \centering
  \includegraphics[width=0.6\linewidth]{ML2Ex1_1_1.jpg}
\end{figure}
\begin{figure}[h]
  \centering
  \includegraphics[width=0.6\linewidth]{ML2Ex1_1_2.jpg}
  \includegraphics[width=0.6\linewidth]{ML2Ex1_1_3.jpg}
\end{figure}

\clearpage

\subsection{Task 1.2: Implementation in code}
\label{subsec:closedform2}

The functions \texttt{add\_bias(X)}, \texttt{fit\_closed\_form(X, y)}, \texttt{predict(X, theta)} and \texttt{mse(y\_true, y\_pred)} and 
\texttt{fit\_closed\_form(X, y)} were implemented.
\newline The linear system $\XX^T\XX\,\tth = \XX^T\yy$ was solved with \texttt{np.linalg.solve} and the results are compared to the manual computations in Task 1.1 (they do match!):
\begin{lstlisting}
=== Task 1.2: Print theta*, the predictions y_hat and the loss function and compare with np.linalg.lstsq ===
theta*: shape = (2,), values = [0.5 0.6]
theta* matches np.linalg.lstsq: True
prediction: shape = (4,), values = [1.1 1.7 2.3 2.9]
prediction matches np.linalg.lstsq: True
training MSE: 0.05
training MSE matches np.linalg.lstsq: True
\end{lstlisting}

\noindent Figure~\ref{fig:toydatafit} plots the toy dataset and displays the linear predictor function \(f(x)\).
\newline As \(f(x)\) is a linear function with parameters $\boldsymbol{\theta} \in \mathbb{R}^2$, $\theta_0^*$ [100 k\$] represents the intercept (value at \(x=0\)), and $\theta_1^*$ [100 k\$\ (per room)] the slope of the straight line.
Parameter $\theta_0^*$ is the bias or in the case of the toy dataset the baseline price for a house (e.g. land price etc., independent of the number of rooms) and is essential to get a proper prediction for the house price (minimum loss).

\begin{figure}[h]
  \centering
  \includegraphics[width=0.6\linewidth]{../Plots/1_toy_fit.pdf}
  \caption{Datapoints of toy dataset with predictions and residuals}
  \label{fig:toydatafit}
\end{figure}

\subsection{Task 1.3: Additional feature: number of windows}
The toy data with the additional window feature from Task1.1(b) was also fitted in the code:
\begin{lstlisting}
Task 1.3: design matrix of shape (4, 3), rank 2, condition number of X^T X: 3.11e+33
fit_closed_form failed: Singular matrix
np.linalg.lstsq: theta* = [0.5  0.12 0.24] , training MSE = 0.0500
\end{lstlisting}

\noindent The output \texttt{fit\_closed\_form failed: Singular matrix} indicates that the matrix
$\mathbf{X}^{\mathsf{T}}\mathbf{X}$ is singular. Since $x_2 = 2x_1$, the corresponding
feature columns in the design matrix $\mathbf{X}$ are linearly dependent, as shown in Task 1.1(b). Therefore,
$\mathbf{X}^{\mathsf{T}}\mathbf{X}$ is non-invertible, and the closed-form solution
cannot be computed. This can be addressed by removing
one of the dependent features - which was also shown in Task 1.1(b): $\boldsymbol{\theta}^1$, $\boldsymbol{\theta}^2$ and $\boldsymbol{\theta}^3$ 
result in the same mean square error - or use more advanced concepts like regularization.

\clearpage

\section{Task 2: Linear Regression on real data}
\label{sec:LRreal}

The provided dataset (\emph{Ames Housing}, $N = 2413$) was imported and randomly split into 70\% training and 30\% test data:
\begin{lstlisting}
========== Task 2 ==========
Ames Housing prices dataset: min = 35.0 k$, median = 159.0 k$, mean = 175.6 k$, max = 755.0 k$
1689 training and 724 test houses - test split = 30.0% of dataset
\end{lstlisting}

\subsection{Task 2.1: Visualisation of dataset}
\label{subsec:LRreal1}

A histogram of the house prices of the dataset is plotted in Figure~\ref{fig:LRReal_histo} and scatter plots of the house price over earch of the 9 included features
is visualized in Figure~\ref{fig:LRReal_poverf}.
\\ The price distribution does \textbf{not appear symmetric}. Reasoning:
\begin{itemize}[itemsep=0.1em]
  \item For a symmetric distribution the left and right sides would be approximately mirror images
  \item The shape of the histogram appear to be right(positively)-skewed with a long right tail with a max\_price = 755.0 k\$
  \item The mean\_price = 175.6 k\$ is significantly larger than the median\_price = 159.0 k\$
\end{itemize}

\noindent Analysing the impact of individual features in relation to the house price in Figure~\ref{fig:LRReal_poverf} leads to following observations:
\begin{itemize}[itemsep=0.1em]
  \item The features with the highest impact on price are:
  \begin{itemize}[itemsep=0.1em]
    \item $Overall$ $quality$: highest impact, price increases almost monotonically as quality increases
    \item $Living$ $area$: clear positive trend
    \item $Garage$ and $basement$ $area$: similiar trends to $living$ $area$ (assumption: correlation with living area - bigger house, bigger basement/garage).
  \end{itemize}
  \item The $overall$ $quality$ relationship to price is clearly non-linear as the higher end of the spectrum increases the price unproportional (non-linear)
  \item The $year$ $built$ relationship seems to be non-linear as well (similiar observations as stated for $overall$ $quality$)
  \item The price of houses realted to the $number$ $of$ $fireplaces$ do not follow a linear trend (e.g. houses with 2 or more fireplaces are not more expensive than houses with less by default)
  \item $Basement$ and $garage$ $area$ do not fit easily in the regression model, as there are several houses with 0$m^2$ in the dataset that span from ~35 k\$ to ~290 k\$.
  \item Besides the houses with 0$m^2$ basements/garages, several other outliers can be detected:
  \begin{itemize}[itemsep=0.1em]
    \item There are several houses with very large lot areas $>10000 m^2$ while most of the houses have $<3000 m^2$
    \item A handful of houses above 600-750 k\$ are extreme price outliers (high overall quality, newly built)
  \end{itemize}
\end{itemize}

\begin{figure}[h]
  \centering
  \includegraphics[width=0.8\linewidth]{../Plots/2_price_histogram.pdf}
  \caption{Histogram of house prices}
  \label{fig:LRReal_histo}
\end{figure}

\begin{figure}[t]
  \centering
  \includegraphics[width=1\linewidth]{../Plots/2_price_over_features.pdf}
  \caption{House price over features}
  \label{fig:LRReal_poverf}
\end{figure}

\clearpage
\subsection{Task 2.2: Linear regression for housing price dataset}
\label{subsec:LRreal2}
The linear regression model from Task 1.2 was fitted for the housing price dataset. Two variantions of the model were implemented:
\begin{itemize}[itemsep=0.1em]
  \item (i):  Model predicts price on the $living$ $area$ only
  \item (ii): Model predicts price on all 9 features of the dataset
\end{itemize}

\noindent The code was adjusted to report $\tth^*$ and the training and test MSE and RMSE $= \sqrt{\text{MSE}}$ of both models:

\begin{lstlisting}
=== Task 2.2: apply linear regression to housing price data ===
Model (i): theta* =
  bias                 16.413
  living_area           1.152
matches np.linalg.lstsq: True

Training set: MSE = 2269.6931 k$^2, RMSE: 47.6413 k$
Training set: MSE matches np.linalg.lstsq: True
Test set: MSE = 2400.8332 k$^2, RMSE: 48.9983 k$

Model (ii): theta* =
  bias               -792.122
  living_area           0.758
  lot_area              0.009
  basement_area         0.326
  garage_area           0.377
  year_built            0.380
  overall_quality      16.245
  bedrooms            -10.539
  bathrooms            -8.145
  fireplaces            6.171
matches np.linalg.lstsq: True

Training set: MSE = 744.1591 k$^2, RMSE: 27.2793 k$
Training set: MSE matches np.linalg.lstsq: True
Test set: MSE = 745.0617 k$^2, RMSE: 27.2958 k$
\end{lstlisting}

\vspace{0.5em}
\noindent $\theta_1^*$ of model (i) in units of [k\$\ (per $m^2$ living area)] can be interpreted as the slope of the linear predictor function 
or in the specific case for the housing price: the predicted price increase in k\$ for increasing the living area by 1 $m^2$.
\vspace{1em}
\\
\noindent Model (i) is a special case of model (ii): the all-features model can reproduce the living-area-only model by setting the coefficients 
of the other eight features to zero. 
Since both models are fitted by least squares on the same training data, model (ii) has at least as many possible predictors and can 
choose the same solution as model (i), or find one with lower training error. 
Therefore, the minimum training MSE of model (ii) cannot be larger that the training MSE of model (i)
This guarantee applies only to training error. Adding features can fit the training data more closely but does not guarantee lower test MSE as
the extra features may cause overfitting.
\vspace{1em}
\\
\noindent The negative bedroom coefficient in model (ii) does not mean that adding a bedroom generally makes a house cheaper. 
In this model, it indicates that, when all other included features are held constant, one additional bedroom is associated with a predicted price decrease of approximately $\$10539$.
In the specific case for housing prices it can be interpreted as follows: with the same living area, overall quality, ... a higher number of bedrooms indicate smaller room sizes on average, which can have an negative impact on the attractiveness of a house and therefore can lower the value of the house overall.

\clearpage

\noindent Figure~\ref{fig:LRReal_linModels} plots the linear model prediction for the house price in k\$ in relation to the incorporated features.
\\Model (i) only use the living area as its sole feature - while the regression line captures the general trend that larger homes tend to be more expensive, the data points scatter around the line, especially for larger living areas. 
This indicates that living area alone cannot fully explain the variation in housing prices, as properties with similar sizes may have very different market values due to other factors.
\\ Predictions from model (ii), which incorporates all available features, are generally clustered more closely around the perfect prediction line, as seen in the right plot, indicating a better fit and higher predictive accuracy.

\begin{figure}[h]
  \centering
  \includegraphics[width=1\linewidth]{../Plots/2_linear_models.pdf}
  \caption{Visualization of the linear models and predictions}
  \label{fig:LRReal_linModels}
\end{figure}

\subsection{Task 2.3: Non-linear features}
\label{subsec:LRreal3}

For task 2.3 $x$ is repleacd by (non-linear) polynomial features $\bm{\phi}(x) = (1, x, x^2, \dots, x^d)^T$ of degree $d$. 
The function \texttt{poly\_features(x, degree)} was implemented, which returns the design matrix $\bm{\Phi} \in \R^{N \times (d+1)}$ including the constant column.
\vspace{0.5em}
\\With this adjusted design matrix the non-linear relationship between the overall quality and the housing price can be better predicted.
\vspace{0.5em}
\\Increasing the polynomial degree allows the model to capture non-linear pattern in the data as seen in the decreasing training MSE
in Figure~\ref{fig:LRReal_nonlinFeatures} and in the output provided:
\vspace{1em}
\begin{lstlisting}
=== Task 2.3: implement non-linear features (poly) ===
degree 1: training MSE =  1882.7 k$^2, test MSE =  1989.3 k$^2
degree 2: training MSE =  1599.9 k$^2, test MSE =  1753.0 k$^2
degree 3: training MSE =  1543.0 k$^2, test MSE =  1722.0 k$^2
degree 4: training MSE =  1539.6 k$^2, test MSE =  1725.0 k$^2
degree 5: training MSE =  1538.9 k$^2, test MSE =  1723.0 k$^2
\end{lstlisting}

\clearpage
\noindent A lower training MSE does not necessarily indicate a better model, since complex models may begin to fit fluctuations and noise rather than the underlying relationship and result in an overfitted model.
The choice of polynomial degree should therefore be based primarily on the test MSE, which measures how well the model generalizes to unseen data. 
The preferred degree is the one that achieves the lowest test MSE or provides the best trade-off between predictive performance and model complexity.
\vspace{0.5em}
\\In the case of the provided housing dataset a \textbf{polynomial degree of 3} would be well fitted for the model.
\vspace{1em}
\\
\noindent Despite using non-linear features, this approach is still considered
\textbf{linear regression} because the model remains linear in its
parameters (weights), e.g. a polynomial model of degree 3 is
\[
\hat{y} = w_0 + w_1x + w_2x^2 + w_3x^3.
\]
Although the model is non-linear in the input variable \(x\), it is linear
in the coefficients \(w_0, w_1, w_2,\) and \(w_3\), since the prediction
is a linear combination of the transformed features the same linear regression framework and optimization methods can be used.
\begin{figure}[h]
  \centering
  \includegraphics[width=1\linewidth]{../Plots/2_polynomial_fits.pdf}
  \caption{Linear regression with non-linear features}
  \label{fig:LRReal_nonlinFeatures}
\end{figure}


\clearpage
\section{Task 3: Logistic Regression}
\label{sec:LogReg}
\subsection{Task 3.1: Pen and paper computation for logistic regression}
\label{subsec:LogReg1}

For task 3.1 several pen and paper computation for the logistic regression were performed:
\vspace{0.5em}
\begin{enumerate}[label=(\alph*), itemsep=0.1em, topsep=0.2em]
    \item Show that $\sigm'(z) = \sigm(z)\big(1 - \sigm(z)\big)$.
    \item For a toy set of $N = 4$ points with $D = 1$: Starting from $\tth^{(0)} = \mathbf{0}$, compute $\mathcal{L}(\tth^{(0)})$, $\nabla_{\tth}\mathcal{L}(\tth^{(0)})$ and $\tth^{(1)}$ for $\eta = 1$. Verify that $\mathcal{L}(\tth^{(1)}) < \mathcal{L}(\tth^{(0)})$.
    \item Show that $\hat{y} \geq 0.5$ if and only if $\tth^T\xx \geq 0$. For $D = 2$ and $\theta_2 \neq 0$, write the decision boundary as $x_2 = g(x_1)$.
\end{enumerate}
\vspace{1em}
\begin{figure}[h]
  \centering
  \includegraphics[width=0.75\linewidth]{ML2Ex3_1_1.jpg}
\end{figure}

\begin{figure}[h]
  \centering
  \includegraphics[width=0.75\linewidth]{ML2Ex3_1_2.jpg}
\end{figure}

\clearpage
\noindent When all parameters of $\tth$ are set to zero, the model assigns probability $\sigma(0) = 0.5$ to each class.
For either label (0 or 1), the binary cross-entropy loss is therefore -log(0.5) = log(2), regardless of how the labels are distributed.
\\ The gradient descent moves closer to the minimum with each iteration step and fits the model better to the training set $\rightarrow$ BCE loss after first iteration is lower than the first guess of $\tth=\mathbf{0}$.

\begin{figure}[h]
  \centering
  \includegraphics[width=0.75\linewidth]{ML2Ex3_1_3.jpg}
\end{figure}

\clearpage
\subsection{Task 3.2 Code implementation of logistic regression}
\label{subsec:LogReg2}

The required functions were implemented in the code skeleton and tested in \texttt{main.py}. 
The pen and paper computations in task 3.1(b) were also validated by fitting the model to the same toy dataset and comparing the results of the output:

\begin{lstlisting}
========== Task 3 ==========
sigmoid(z) = [0.00000000e+00 4.53978687e-05 5.00000000e-01 9.99954602e-01
 1.00000000e+00] (no overflow)
gradient_descent: starting at theta0 = [0. 0.] with learning_rate = 1.0 and eps = 1e-06
Task 3.1(b): L(theta^(0)) = 0.6931, gradient = [ 0.  -0.5], theta^(1) = [0.  0.5], L(theta^(1)) = 0.4886
\end{lstlisting}
\vspace{1em}
\noindent The housing dataset was classified/labeled as "high prices", if the house price exceeds $\geq 250\,\text{k\$}$, the dataset was split 
into the training and test sets and their features $living$ $area$ and $overall$ $quality$ were standardized for further computations.
\vspace{0.5em}
\begin{lstlisting}
fraction of houses with price >= 250 k$: 0.125 (training), 0.149 (test)
standardization of ['living_area', 'overall_quality']: mean [137.07   6.  ], std [44.69  1.33]
\end{lstlisting}
\vspace{1em}
\noindent After running some sanity checks on the implemented functions, the model classifies the houses from their standardized living area and 
overall quality with the gradient descent optimization algorithm ($\eta = 1$, $\varepsilon = 10^{-6}$, at most $5000$ iterations). 
\\ The number of iterations, the final training loss, $\tth$ and the training and test accuracy are reported and
the plots for decision boundary and the loss function over iterations are included:
\vspace{0.5em}
\begin{lstlisting}
L(0) = log 2: True
gradient_bce vs. finite differences, max. abs. difference: 2.28e-10
gradient_descent: starting at theta0 = [0. 0. 0.] with learning_rate = 1.0 and eps = 1e-06

Task 3.2: number of iterations = 2653
Task 3.2: final training loss = 0.1351
Task 3.2: final theta = [-5.3571242   2.01155552  2.84604839]
Task 3.2: training accuracy = 0.9455
Task 3.2: test accuracy = 0.9351
\end{lstlisting}
\vspace{1em}
\noindent Figure~\ref{fig:LogReg_decisionBound} shows the decision boundary of a binary classifier that separates houses priced below \$250k ("not high-priced") from houses priced at or above \$250k ("high-priced")
using $living$ $area$ and $overall$ $quality$ as input features. The data reveal a clear relationship between overall quality and house price: 
houses with higher quality ratings are much more likely to belong to the high-price class. The fitted decision boundary has a negative slope, 
indicating that a smaller living area can still lead to a high predicted price if the overall quality is sufficiently high, 
whereas larger houses with lower quality ratings are more likely to be classified as lower-priced properties.
The dotted lines represent predicted probabilities of 0.1 and 0.9 and illustrate the uncertainty region around the decision boundary.

\vspace{0.5em}
\noindent Figure~\ref{fig:LogReg_lossCurve} illustrates the training loss measured by binary cross-entropy (BCE) as a function of the optimization iterations. 
The loss decreases monotonically and indicates that the optimization algorithm successfully updates the model parameters and progressively 
improves the fit to the training data.
Since the loss stabilizes and no oscillations or divergence are observed, the optimization appears to be numerically 
stable and effective for the given classification task.

\begin{figure}[h]
  \centering
  \includegraphics[width=1\linewidth]{../Plots/3_decision_boundary.pdf}
  \caption{Decision boundary for highly prices houses}
  \label{fig:LogReg_decisionBound}
\end{figure}

\begin{figure}[h]
  \centering
  \includegraphics[width=1\linewidth]{../Plots/3_loss_curve.pdf}
  \caption{Binary cross entropy loss (BCE) over iterations for training set}
  \label{fig:LogReg_lossCurve}
\end{figure}

\clearpage
\noindent When only about 13\% of the houses are high-priced, a classifier that always predicts “not high-priced” gets roughly 
87\% accuracy. In this dataset, the baseline is 1 - 0.125 = 87.5\% on training set and 1 - 0.149 = 85.1\% on the test set.
The model\'s test accuracy of 93.51\% is about 8\% higher than the test baseline which is a clear improvement, however the model may still perform poorly on outliers.

\vspace{1em}

\begin{figure}[h]
  \centering
  \includegraphics[width=0.75\linewidth]{ML2Ex3_3.jpg}
\end{figure}


\noindent With $\tth$ and Task~3.1(c) the minimal (standardized) overall quality for which a house of average living area (standardized value $0$) is classified as high-priced was calculated to $\mathbf{1.822}$.
As the overall quality is defined as a discrete (integer) value, the minimum (non-standardized) overall quality that would label a house of average living area as "high-priced" is $\mathbf{9}$.

\vspace{2em}
\noindent \textit{Additional sanity checks for the last task:}
\vspace{0.5em}
\\
If we would draw a straight vertical line in Figure~\ref{fig:LogReg_decisionBound} for the average living area ($137m^2$), it would hit the decision 
boundary at the calculated non-standardized overall quality around 8.4, which is consistent with the pen and paper computation.
\\ Also a sanity check was done in the python implementation. For the standardized living area = 0, the model should predict "not high priced" for standardized
overall quality <1.882 and "high priced" for values higher than that.
\vspace{0.5em}
\begin{lstlisting}
Sanity check Task 3.2: predict label for average living area and std_quality 1.88: 0
Sanity check Task 3.2: predict label for average living area and std_quality 1.90: 1
\end{lstlisting}

\clearpage
\section{Addendum}
\subsection{Usage of generative AI}
\noindent Generative AI was used in the the form of GitHub Copilot in Visual Studio Code. It was used to debug code, 
especially for pyhton syntax and dimension errors, and to format LaTeX for the report.
Another usecase was to generate the subplots in task 2.1 (visualisation of dataset) as it is was a repetetive task to set all labels, format axis etc.
\vspace{0.5em}
\\The code generation/autocompletion feature was used in some cases, but required manual checking and adjustments along the way 
as the AI model made wrong assumptions for the implementation or wrote code in a hard-to-follow style and tried to put the whole
functionality in a single return line which made it harder to understand or debug the functions.
\vspace{0.5em}
\\Copilot was also used to analyse the chart for the decision boundary to provide some pointers to include in the report, which worked better than 
expected with only using simple open-ended prompts such as "What do you see in this chart?"

\end{document}